% =========================================================
% IV. LEGALSCM-RAG
% =========================================================
\section{LegalSCM-RAG}
\label{sec:legalscm_rag}

Chúng tôi đề xuất \textit{LegalSCM-RAG}, một kiến trúc RAG cho hỏi đáp pháp
luật nhiều bước, kết hợp truy xuất trên đồ thị điều kiện--hệ quả, suy luận theo
đường nhiều bước và xác minh chọn lọc bằng mô hình nhân quả cấu trúc pháp lý
(\textit{Legal Structural Causal Model -- LegalSCM}).

Với câu hỏi $q$, pipeline tổng thể là
\begin{equation}
\begin{aligned}
q
&\rightarrow
\text{Causal Retrieval}
\rightarrow
\mathcal{P}(q)
\rightarrow
P^*\\
&\rightarrow
\text{Question-Aware Verification}
\rightarrow
\mathcal{E}^*
\rightarrow
(a,C),
\end{aligned}
\label{eq:overall_pipeline}
\end{equation}
trong đó $P^*$ là đường suy luận chính, $\mathcal{E}^*$ là tập bằng chứng được
chọn, $a$ là câu trả lời và $C$ là tập căn cứ pháp lý.

\subsection{Truy xuất dựa trên đồ thị pháp luật}

Hệ thống biểu diễn câu hỏi bằng embedding và truy xuất các sự kiện, quy tắc
liên quan từ causal memory. Các sự kiện truy xuất đóng vai trò seed để mở rộng
trên đồ thị. Từ mỗi seed, hệ thống duyệt các quan hệ điều kiện--hệ quả có hướng
để xây dựng
\begin{equation}
\mathcal{P}(q)=\{P_1,P_2,\ldots,P_K\},
\end{equation}
trong đó
\begin{equation}
P_i=v_0\xrightarrow{r_1}v_1\xrightarrow{r_2}\cdots
\xrightarrow{r_h}v_h.
\end{equation}
Khác RAG ngữ nghĩa thông thường, đơn vị truy xuất bao gồm cả sự kiện pháp lý và
các quan hệ quy tắc có hướng nối chúng, thay vì chỉ gồm các đoạn văn độc lập.

\subsection{Xếp hạng đường suy luận}

Các đường ứng viên được xếp hạng bằng tín hiệu ngữ nghĩa và cấu trúc:
\begin{equation}
S(P_i,q)=
\alpha S_{\mathrm{sem}}+
\beta S_{\mathrm{graph}}+
\gamma S_{\mathrm{seed}}+
\delta S_{\mathrm{path}}.
\end{equation}
Đường chính được chọn bởi
\begin{equation}
P^*=\arg\max_{P_i\in\mathcal{P}(q)}S(P_i,q).
\end{equation}
Các quy tắc và điều khoản gắn với $P^*$ được chuyển tới tầng xác minh dưới dạng
bằng chứng có cấu trúc.

\subsection{Suy luận điểm bất động LegalSCM}

LegalSCM là mô hình quy tắc quy phạm ba giá trị và có tính xác định. Mỗi sự kiện
thuộc $\{\mathrm{TRUE},\mathrm{FALSE},\mathrm{UNKNOWN}\}$. Một quy tắc có
dạng
\begin{equation}
X_1\land\cdots\land X_k\rightarrow Y.
\end{equation}
Quy tắc chỉ kích hoạt khi mọi literal điều kiện có trạng thái xác định đúng yêu
cầu và không có ngoại lệ được mã hóa nào đang hoạt động. Vì vậy, các điều kiện
trong một quy tắc có quan hệ hội. Nhiều quy tắc kích hoạt độc lập cùng suy ra
một trạng thái của $Y$ biểu diễn hỗ trợ thay thế theo quan hệ tuyển. Các hỗ trợ
trái dấu được giữ như xung đột và quy về trạng thái hiệu dụng
\textsc{unknown}. Toán tử quy tắc được lặp tới điểm bất động như trong
Mục~\ref{sec:problem_formulation}.

Engine chuẩn hỗ trợ nhiều điều kiện và ngoại lệ, nhưng dữ liệu legacy trong
thực nghiệm hiện tại chỉ tạo một điều kiện dương và một hệ quả dương cho mỗi
quy tắc. Do đó, kết quả hiện tại chưa phải là đánh giá thực nghiệm cho
conjunction nhiều điều kiện hoặc xử lý ngoại lệ tường minh.

LegalSCM mô hình hóa các phụ thuộc quy phạm trích xuất từ luật thành văn; mô
hình không học quan hệ nhân quả từ dữ liệu quan sát và không ước lượng hiệu ứng
nhân quả thống kê.

\subsection{Hard intervention và topological diagnostic}

Với câu hỏi yêu cầu rõ ràng việc loại bỏ sự kiện trung gian $M$, LegalSCM thực
hiện hard intervention tương đối với mô hình quy tắc đã mã hóa
\cite{pearl2009causality}:
\begin{equation}
do(X_M=0),\qquad 0\equiv\mathrm{FALSE}.
\end{equation}
Phép can thiệp cố định $X_M$ ở false, chặn mọi cơ chế quy tắc có hệ quả là
$X_M$, rồi suy luận lại điểm bất động. Nếu đích factual $Y$ là true và
\begin{equation}
Y_{do(X_M=0)}=\mathrm{FALSE},
\end{equation}
thì $M$ là cần thiết trong ngữ cảnh đã mã hóa. Nếu
$Y_{do(X_M=0)}=\mathrm{TRUE}$, một cơ chế thay thế vẫn hỗ trợ đích. Nếu một
trong hai thế giới là unknown, kết quả chưa xác định.

Phép toán này khác với graph ablation phụ trợ tạo
$G^{-M}=G[V\setminus\{M\}]$ rồi tìm đường thay thế có giới hạn. Phép ablation
xóa nút cùng các cạnh kề và chỉ được giữ như diagnostic/fallback về topology;
nó không phải Pearl-style intervention. Ngoài ra, LegalSCM không thực hiện
abduction trạng thái ngoại sinh. Kết quả vì vậy là tính cần thiết tương đối với
mô hình quy phạm và context đầu vào, không phải actual causation ở cấp cá thể.

\subsection{Xác minh phụ thuộc câu hỏi}

Verifier chọn một trong ba route từ nội dung câu hỏi. Một claim trực tiếp
$A\rightarrow C$ được kiểm tra trên hai endpoint được ground thận trọng trong
đồ thị. Câu hỏi pháp lý thông thường nhận factual validation của Primary Path
bằng LegalSCM mà không can thiệp. Chỉ câu hỏi nêu rõ việc loại mediator hoặc
tính cần thiết mới yêu cầu $do(X_M=0)$ và so sánh hai thế giới. Routing này
ngăn direct-edge topology hoặc factual validation bị gọi nhầm là structural
intervention.

\subsection{Cơ chế xác minh chọn lọc}

Hệ thống phân biệt miền trạng thái sự kiện
$\{\mathrm{TRUE},\mathrm{FALSE},\mathrm{UNKNOWN}\}$, các trạng thái can thiệp
nội bộ như necessary/non-necessary, và miền nhãn cuối
\begin{equation}
\mathcal{Z}=\{\textsc{supported},\textsc{reject-direct-claim},
\textsc{uncertain}\}.
\end{equation}
Nhãn verifier chỉ có tính ràng buộc khi vượt qua các điều kiện về confidence,
endpoint hoặc mediator grounding, độ phủ Primary Path và tính nhất quán trạng
thái:
\begin{equation}
z=
\begin{cases}
z_{\mathrm{SCM}}, & G_{\mathrm{commit}}=1,\\[3pt]
z_{\mathrm{LLM}}, & G_{\mathrm{commit}}=0.
\end{cases}
\end{equation}
Khi gate thất bại, verifier không ép nhãn \textsc{uncertain}. Thay vào đó, bằng
chứng Causal Path RAG không đổi được chuyển cho mô hình sinh chung. Full
LegalSCM do đó là một hệ lai chọn lọc, không phải classifier do SCM quyết định
trên mọi mẫu.

\subsection{Tinh lọc bằng chứng và sinh câu trả lời}

Sau xác minh, hệ thống tạo
\begin{equation}
\mathcal{E}^*=\{r_j\mid r_j\text{ hỗ trợ }P^*\}
\end{equation}
và cung cấp
\begin{equation}
I_q=(q,P^*,z,\mathcal{E}^*)
\end{equation}
cho mô hình ngôn ngữ để sinh
\begin{equation}
(a,C)=f_{\theta}(I_q).
\end{equation}
Việc tách các tầng giúp truy xuất, suy luận quy tắc, route xác minh và sinh ngôn
ngữ có thể được kiểm tra độc lập.
